% 18 November 2026. Drafted for Edward Baker with OpenAI Codex assistance.
% Actual model and effort unavailable; helper settings inherited.
\section{A joint comparison in physical backward time}
\label{sec:jointtime}

The initial-environment comparison discards the positive forcing in the
equation for $M$. This can lose a cancellation that persists throughout
the evolution. Section~\ref{sec:canonicaltime} will show that, at high
canonical zeros, even exact initial interaction data give that comparison
only an $O(g^2)$ duration. We first develop a different interface: a lower
bound for a combination of the evolving gap and its complete exterior
moments. The comparison and its finite-head acquisition are from
\cite{BakerCM23,BakerCM25}; the role of the source, rather than additional
initial data alone, will be explicit.

\subsection{The defect equation and its barrier}

Use the physical backward variable $u=t_*-t$, and initially work in the
stopped all-real simple regime of Section~\ref{sec:statistics}. For one
adjacent pair define the physical forcing and two dimensionless quantities
\begin{equation}
 \mathcal F=12S_2+16DS_3,\qquad
 V=1-2DM,\qquad
 K=2D^2\mathcal F=24D^2S_2+32D^3S_3.
 \label{jt:variables}
\end{equation}
Thus $\mathcal F$ differs by a factor $\ell^{-4}$ from the normalized
forcing $F$ of Section~\ref{sec:joint}. All sums in
\eqref{jt:variables} concern the same pair and the complete divisor.
There is no moment cutoff in the following identities.

\begin{proposition}[Joint defect equation]
\label{jt:defect}
Where $D>0$ and the complete pair laws hold,
\begin{equation}
 D_u=-2V,\qquad V_u=\frac{1-V^2-K}{D}.
 \label{jt:system}
\end{equation}
\end{proposition}

\begin{proof}
The first identity is the pair equation. Using
$M_u=-2M^2+\mathcal F$ in the derivative of $V$ gives
\[
 V_u=-2(D_uM+DM_u)=4M-4DM^2-2D\mathcal F.
\]
Since $2DM=1-V$, the first two terms equal $(1-V^2)/D$,
and the last equals $-K/D$.
\end{proof}

\begin{theorem}[Joint forcing barrier]
\label{jt:barrier}
Suppose the real complete pair laws hold on a stopped interval
$0\le u<\tau$ with $D>0$, and suppose
\begin{equation}
 K(u)\ge1-\eta^2,\qquad V(0)\le v_0,
 \qquad \eta\ge0.
 \label{jt:barrier-input}
\end{equation}
For any $v\ge\max\{v_0,\eta\}$,
\begin{equation}
 V(u)\le v,\qquad D(u)\ge D(0)-2vu.
 \label{jt:barrier-output}
\end{equation}
In particular, if $D(0)\ge d_0>0$ and
\begin{equation}
 2v\Delta<d_0,
 \label{jt:strict-margin}
\end{equation}
the selected pair cannot collide by $u=\Delta$ while the stipulated
regime and inputs remain available up to a first event. When $v=0$,
$D$ is nondecreasing on that regime.
\end{theorem}

\begin{proof}
On a compact subinterval before the stop, $D$ has a positive minimum.
For $w=V-v$, \eqref{jt:system} gives
\[
 w_u+\frac{V+v}{D}w
 \le\frac{\eta^2-v^2}{D}\le0.
\]
The coefficient is continuous. Its integrating factor is positive,
regardless of the sign of $V+v$, and $w(0)\le0$. Therefore $w\le0$.
Integration of $D_u=-2V$ proves \eqref{jt:barrier-output}.
Exhausting the stopped interval by compact subintervals preserves the
bound. At a putative collision time $u_c\le\Delta$, continuity of the
gap gives $D(u)\to0$, contradicting
$D(u)\ge d_0-2v\Delta>0$. No division by $D$ is made at the event.
\end{proof}

For a gap with $D(0)\asymp L^{-2}$, sufficient rates for this comparison
to pay a fixed available duration are
\begin{equation}
 V(0)_+=o(L^{-2}),\qquad
 1-K(u)\le\eta^2=o(L^{-4})\quad\hbox{uniformly in }u.
 \label{jt:rates}
\end{equation}
The square root of the forcing deficiency is essential. A statement
$K=1+o(1)$ alone does not give \eqref{jt:rates}. These are sufficient
rates for the comparison, not necessary conditions for actual
noncollision. In the global all-real application, a different earlier
unprotected collision ends the sign regime; the theorem does not
continue the argument past it.

\begin{remark}[Local signed complete version]
\label{jt:local-version}
The proof of Theorem~\ref{jt:barrier} uses neither $M\ge0$ nor positivity
of the individual exterior summands. It therefore applies to a
conjugation-invariant count-two tube with a possibly nonreal exterior,
provided its complete product, contour and first-time-jet passages have
been paid on the stated time domain. To see the identity transfer,
use the fixed cleared contour and the complete product as in
Lemma~\ref{lem:slice-passage}. The finite polynomial identities for $D_u$
and $M_u$ pass through uniform spatial jets on that contour and the
absolute inverse-square tails. Reality follows from conjugation.
Before the first local loss of $D>0$, the same integrating-factor proof
then applies to a directly acquired signed complete bound for $K$.
Positive omission of an exterior tail is not available in this version.
Section~\ref{sec:compact} formulates its source acquisition on a stopped
trial-coordinate domain, without assuming positivity on the desired
whole interval. At a regular target merger $\mathcal F$ is finite and
$K=2D^2\mathcal F=0$; thus an unconditional positive floor for $K$ on
the whole tube would itself already exclude that merger.
\end{remark}

\subsection{Finite neighboring heads and an explicit perturbation bound}

Consider first an equally spaced real divisor, a target pair $0,g$, and
$J$ exterior zeros on each side. Put $D=g^2/4$. The $m$th exterior
occurrence on either side has denominator $Q=m(m+1)g^2$. Denote the
interaction subsum by $b_J$ and the forcing subsum by $K_J$.

\begin{lemma}[Exact lattice head]
\label{jt:lattice}
For every integer $J\ge1$,
\begin{equation}
 2Db_J=\frac{J}{J+1},\qquad K_J=1-\frac1{(J+1)^3}.
 \label{jt:lattice-values}
\end{equation}
\end{lemma}

\begin{proof}
The two occurrences at level $m$ contribute
$2/[g^2m(m+1)]$ to $b_J$, so the first equality follows from
$1/[m(m+1)]=1/m-1/(m+1)$. Their combined contribution to $K_J$ is
\[
 \frac3{m^2(m+1)^2}+\frac1{m^3(m+1)^3}
 =\frac1{m^3}-\frac1{(m+1)^3},
\]
because $3m(m+1)+1=(m+1)^3-m^3$. Summing proves the second equality.
\end{proof}

For an all-real adjacent pair, omitted terms of $K$ and $M$ are
nonnegative. A finite lower bound for $K$ therefore needs no subtraction
of an unknown upper tail. This differs from a complete asymptotic
equality for $M$, whose remaining tail must be bounded. The following
version keeps explicit constants for a perturbed head.

\begin{proposition}[Robust finite-head bounds]
\label{jt:robust-head}
At a real-simple slice, let the target gap be $g>0$, let $d>0$, and
let $0\le e<1/3$. Suppose $g\ge(1-e)d$. For the $m$th selected
exterior occurrence on each side, $1\le m\le J$, let $r$ be its
distance to the nearer endpoint, and suppose
\begin{equation}
 r\le(1+2e)md,\qquad r+g\le(1+2e)(m+1)d.
 \label{jt:head-distances}
\end{equation}
With $b_J$ the interaction of these selected occurrences,
\begin{equation}
 \begin{split}
 K&\ge1-18e-(J+1)^{-3},\\
 1-2DM&\le1-2Db_J\le6e+(J+1)^{-1}.
 \end{split}
 \label{jt:robust-bounds}
\end{equation}
If the distance hypotheses hold with one $e_*$ throughout a stopped
all-real interval, an admissible barrier parameter is
\begin{equation}
 v=\max\left\{6e_*+(J+1)^{-1},
              \sqrt{18e_*+(J+1)^{-3}}\right\}.
 \label{jt:head-barrier}
\end{equation}
Only the initial instance of the second inequality is needed here.
\end{proposition}

\begin{proof}
Put $\gamma=(1-e)/(1+2e)\in[0,1]$. Since $Q=r(r+g)$,
\[
 \frac{D}{Q}\ge\frac{\gamma^2}{4m(m+1)}.
\]
The quadratic and cubic terms in $D/Q$ contributing to $K$ are at
least $\gamma^4$ and $\gamma^6$ times the lattice terms. Both are at
least $\gamma^6$ times those terms. Positivity of the omitted complete
exterior and Lemma~\ref{jt:lattice} give
\[
 K\ge\gamma^6[1-(J+1)^{-3}],\qquad
 2DM\ge2Db_J\ge\gamma^2\frac{J}{J+1}.
\]
Now $1-\gamma\le3e$, so
$1-\gamma^6\le18e$ and $1-\gamma^2\le6e$. This proves
\eqref{jt:robust-bounds}. Insert its uniform forcing bound and its
initial defect bound into Theorem~\ref{jt:barrier} to obtain
\eqref{jt:head-barrier}.
\end{proof}

With exact lattice head inequalities, $e_*=0$ and
$v=(J+1)^{-1}$. The strict duration becomes
\begin{equation}
 \Delta<\frac{g(0)^2(J+1)}8.
 \label{jt:head-duration}
\end{equation}
The growing allowance uses control of the head throughout time. Adding
the same neighbors only to the initial mass does not supply it. Nor is
this a finite divisor evolution: the selected labels and $D$ are
evaluated in the unchanged complete heat flow.

\subsection{A ratio-forcing floor retains a different loss}

To distinguish the joint quantity from a weaker datum, suppose $M>0$
and write $\kappa=\mathcal F/(2M^2)$. Equation~\eqref{jt:system} becomes
\begin{equation}
 V_u=2M\{(1-\kappa)+(1+\kappa)V\},
 \qquad K=(1-V)^2\kappa.
 \label{jt:ratio-system}
\end{equation}
Even near $\kappa=1$, an upper comparison for a positive defect has
coefficient approximately $4M$. The following exact integration records
the resulting allowance, rather than asserting instability of the
underlying heat flow.

\begin{proposition}[Ratio-floor comparison]
\label{jt:ratio-comparison}
In the stopped all-real regime suppose
$\mathcal F\ge2(1-\varepsilon)M^2$, where $0<\varepsilon<1$,
and put $b=M(0)>0$ and $A(u)=1+2\varepsilon bu$. Then
\begin{equation}
 \begin{split}
 M(u)&\ge b/A(u),\\
 D(u)&\ge A(u)^{2/\varepsilon}
 \left[D(0)-
 \frac{1-A(u)^{-(2-\varepsilon)/\varepsilon}}
      {b(2-\varepsilon)}\right].
 \end{split}
 \label{jt:ratio-bound}
\end{equation}
If $r=1-(2-\varepsilon)bD(0)>0$, positivity of this lower bound
holds precisely while
\begin{equation}
 u<\frac{r^{-\varepsilon/(2-\varepsilon)}-1}
          {2\varepsilon b}.
 \label{jt:ratio-horizon}
\end{equation}
For $\varepsilon=0$ the corresponding lower bound is
\begin{equation}
 D(u)\ge\e^{4bu}
       \left[D(0)-\frac{1-\e^{-4bu}}{2b}\right].
 \label{jt:ratio-zero}
\end{equation}
\end{proposition}

\begin{proof}
The hypothesis gives $M_u\ge-2\varepsilon M^2$. The same
integrating-factor comparison as in
Theorem~\ref{stat:backward-comparison} yields $M\ge b/A$.
Before collision $D\ge0$, so multiplying
$D_u\ge-2+4Db/A$ by $A^{-2/\varepsilon}$ and integrating gives
\eqref{jt:ratio-bound}; explicitly,
\[
 2\int_0^u A(v)^{-2/\varepsilon}\dd v
 =\frac{1-A(u)^{-(2-\varepsilon)/\varepsilon}}
       {b(2-\varepsilon)}.
\]
Solving the strict bracket inequality proves
\eqref{jt:ratio-horizon}. When $\varepsilon=0$, $M\ge b$ and
the integrating factor is $\e^{-4bu}$, giving
\eqref{jt:ratio-zero} directly.
\end{proof}

For a positive initial defect $\delta=1-2bD(0)$, the last bound has
horizon $\log(1/\delta)/(4b)$. At $bD(0)=1/2$,
\eqref{jt:ratio-horizon} is asymptotic to
$\log(2/\varepsilon)/(4b)$ as $\varepsilon\downarrow0$.
Thus with $b\asymp L^2$, a conservative positive defect allowance of
size $L^{-1}$ gives a duration of order $(\log L)/L^2$; an allowance
of size $\e^{-cL}$ gives order $L^{-1}$. This particular absolute-defect
comparison would require accuracy of order $\e^{-c_\Delta L^2}$
to pay a fixed duration. A favorable proved sign or additional
correlation is outside this limitation. No positive sign of the actual
canonical defect, or lower bound on its actual error, is asserted.

An exact heat control separates initial near-spacing from a whole-time
forcing floor. For $k>0$ and $0<\alpha_0<1$, consider
\begin{equation}
 H(u,z)=\e^{-k^2u}\cos(kz)+\alpha_0,
 \qquad H_u=H_{zz}.
 \label{jt:cosine-control}
\end{equation}
This is an even finite positive atomic source, not theta. Initially all
zeros are real and simple. Put $\alpha(u)=\alpha_0\e^{k^2u}$ and
$\theta(u)=\arccos\alpha(u)$. The short pair around $\pi/k$ has
\[
 D(u)=\theta(u)^2/k^2,\qquad
 V(u)=\frac{\theta(u)\alpha(u)}{\sqrt{1-\alpha(u)^2}},
 \qquad u_c=k^{-2}\log(1/\alpha_0).
\]
These formulas follow by solving $\cos(kz)=-\alpha(u)$ and
differentiating; in particular $D_u=-2V$. As $\alpha_0\downarrow0$,
the initial gaps differ from $\pi/k$ by $O(\alpha_0/k)$ and
$V(0)=(\pi/2)\alpha_0+O(\alpha_0^2)$. Choosing
$\alpha_0=\e^{-ck}$ makes those errors exponentially small, yet
$u_c=c/k$. For each fixed $k$, the other clusters stay separated from
the selected pair at $u_c$, and their periodic inverse-power tails
converge uniformly. Hence $S_2,S_3$ stay bounded while $D\to0$, giving
$K\to0$. The near-unit whole-time forcing premise fails.

The stationary cosine heat field has $V=0$, $K=1$ and is preserved
exactly by the joint barrier. The quadratic collision field has
$M=K=0$, $V=1$ and requires $\eta\ge1$, reproducing its gap-only
lifetime. These controls satisfy the local equations but lack the
unchanged canonical source and population. They test the hypotheses
and strict endpoint interpretation, not the source-specific acquisition
proved next.
